\chapter{Transformées, interpolation et différentiation algorithmique}\label{ch:qm:transforms-interpolation-and-algorithmic-differentiation}

Le calcul de risque déplace 400 points de la courbe un à un et revalorise le portefeuille 401 fois chaque nuit. Les mêmes 400 sensibilités sortent d'un seul enregistrement du code de valorisation et d'un seul parcours
à rebours de celui-ci, pour le coût de quelques valorisations. Ce chapitre rassemble les outils numériques dont se composent les codes de valorisation et de risque~: méthodes de Fourier qui valorisent à partir d'une \omterm{def:qm:probability-at-speed:charfn}{fonction caractéristique}, interpolation qui
transforme une poignée de cotations en courbe, méthodes de recherche de racines qui inversent une formule de valorisation, et \omterm{def:qm:transforms-interpolation-and-algorithmic-differentiation:ad}{différentiation algorithmique}, qui fournit les dérivées exactes de tout ce que calcule le code.

\section{Méthodes de Fourier}

\begin{definition}[Transformées de Fourier discrète et rapide]\label{def:qm:transforms-interpolation-and-algorithmic-differentiation:dft}
La \emph{transformée de Fourier discrète}\index{transformée de Fourier discrète} de $x_0, \dots, x_{N-1}$ est $X_k = \sum_{j=0}^{N-1}x_je^{-2\pi\mathrm ijk/N}$, $k = 0, \dots, N - 1$. Une \emph{transformée de Fourier
rapide}\index{transformée de Fourier rapide} la calcule en $O(N\log N)$ opérations au lieu de $O(N^2)$ en la découpant récursivement en transformées de moitié moindre longueur (Cooley et Tukey, 1965).
\end{definition}

Un modèle de Lévy est spécifié par son \omterm{def:qm:jump-processes:triplet}{exposant caractéristique} (chapitre~6), et non par sa densité, de sorte que les prix découlent de la \omterm{def:qm:probability-at-speed:charfn}{fonction caractéristique} $\varphi(u) = \E[e^{\mathrm iu\ln(S_T/S_0)}]$. Carr et Madan (1999)
valorisent d'un coup toute une gamme de strikes à l'aide d'une seule \omterm{def:qm:transforms-interpolation-and-algorithmic-differentiation:dft}{transformée de Fourier rapide} d'un call amorti~; la méthode ci-dessous ne nécessite aucun amortissement et converge plus vite.

\begin{proposition}[Inversion de Fourier]\label{prop:qm:transforms-interpolation-and-algorithmic-differentiation:gp}
Pour une variable aléatoire $X$ de \omterm{def:qm:probability-at-speed:charfn}{fonction caractéristique} $\varphi$ et un point de continuité $k$ de sa loi, $\P(X > k) = \frac12 + \frac1\pi\int_0^\infty\mathrm{Re}\bigl[e^{-\mathrm iuk}\varphi(u)/(\mathrm iu)\bigr]\,du$
(Gil-Pelaez, 1951).
\end{proposition}

Pour une \omterm{def:qm:jump-processes:jd}{diffusion à sauts} de Merton ($\sigma = 15\%$, sauts à l'intensité 0.5 par an d'espérance $-0.1$ et d'écart type 0.2 en log-prix, $r = 3\%$, un an), la formule avec la règle des trapèzes sur $[0, 200]$ donne
$\P(S_T > 110) = 0.3503999$, à comparer au mélange de Poisson de lognormales, à $3 \times 10^{-10}$ près.

\begin{definition}[Méthode COS]\label{def:qm:transforms-interpolation-and-algorithmic-differentiation:cos}
La \emph{méthode COS}\index{méthode COS} (Fang et Oosterlee, 2008) tronque la densité de $y = \ln(S_T/K)$ à un intervalle $[a, b]$, la développe en série de cosinus dont les coefficients se lisent sur la
\omterm{def:qm:probability-at-speed:charfn}{fonction caractéristique}, $A_k \approx \frac2{b-a}\mathrm{Re}[\varphi(\frac{k\pi}{b-a})e^{-\mathrm ik\pi a/(b-a)}]$, et intègre le payoff contre chaque cosinus sous forme fermée, de sorte qu'un prix soit
$e^{-rT}\sum_{k<N}{}'A_kV_k$ (le premier terme étant divisé par deux).
\end{definition}

\begin{proposition}[Convergence de la méthode COS]\label{prop:qm:transforms-interpolation-and-algorithmic-differentiation:cosconv}
Pour une densité régulière sur l'intervalle de troncature, les coefficients en cosinus décroissent plus vite que toute puissance de $k$, et l'erreur du prix COS décroît exponentiellement en $N$ jusqu'à ce que l'\omterm{def:qm:finite-difference-methods:consistent}{erreur de troncature}
de $[a, b]$ domine (Fang et Oosterlee, 2008~; résultat admis).
\end{proposition}

Avec $[a, b]$ fixé à dix écarts types à partir des \omterm{def:qm:jump-processes:cumulant}{cumulants}, un call de Black--Scholes est exact à $7 \times 10^{-15}$ près avec 64 termes, et le call de Merton à $10^{-7}$ près avec 64 termes et à $2 \times 10^{-12}$ près avec 128
(voir la \cref{fig:qm:transforms-interpolation-and-algorithmic-differentiation:cos}). La maison valorise les calls à partir du développement des puts et de la parité, ce qui maintient les coefficients bornés. Pour un modèle variance-gamma,
où seule la simulation sert de référence, 256 termes donnent 9.5243 et 400\,000 trajectoires simulées $9.528 \pm 0.022$.

\begin{omfigure}
\begin{tikzpicture}
\begin{semilogyaxis}[width=10cm,height=5cm,xlabel={nombre de termes en cosinus $N$},ylabel={erreur de prix},
  tick label style={font=\small},label style={font=\small},xmin=0,xmax=264,ymin=1e-15,ymax=10,grid=major,grid style={gray!25},
  legend cell align=left,legend style={font=\footnotesize,at={(0.97,0.97)},anchor=north east,fill=white,draw=none},table/col sep=comma]
\addplot[omDef,very thick,mark=*,mark size=1.5pt] table[x=N,y=bs]{figdata/methods/28-transforms-interpolation-and-algorithmic-differentiation/cos.csv};
\addplot[omThm,very thick,mark=square*,mark size=1.5pt] table[x=N,y=merton]{figdata/methods/28-transforms-interpolation-and-algorithmic-differentiation/cos.csv};
\legend{Black--Scholes ($\sigma = 20\%$),diffusion à sauts de Merton}
\end{semilogyaxis}
\end{tikzpicture}

\omcaption{Erreur du prix COS d'un call à la monnaie d'un an en fonction du nombre de termes en cosinus (axe des erreurs logarithmique, $N$ linéaire)~: convergence exponentielle jusqu'à l'arrondi pour Black--Scholes et
à la précision de la série de référence pour Merton. Données~: le tutoriel du chapitre.}\label{fig:qm:transforms-interpolation-and-algorithmic-differentiation:cos}
\end{omfigure}

\section{Interpolation et splines}

\begin{definition}[Splines cubiques, naturelles et monotones, phénomène de Runge]\label{def:qm:transforms-interpolation-and-algorithmic-differentiation:spline}
Une \emph{spline cubique}\index{spline cubique} passant par $(x_i, y_i)$ est une fonction cubique par morceaux dont les dérivées première et seconde sont continues~; une \emph{spline cubique naturelle}\index{spline cubique naturelle} a une
dérivée seconde nulle aux deux extrémités, et ses dérivées secondes aux nœuds résolvent un système tridiagonal. L'\emph{interpolation cubique monotone}\index{interpolation cubique monotone} choisit les pentes aux nœuds
d'une interpolante de Hermite cubique par morceaux de sorte qu'elle soit monotone partout où les données le sont. Le \emph{phénomène de Runge}\index{phénomène de Runge} (Runge, 1901) est la croissance, près des extrémités de l'intervalle, de l'erreur de
l'interpolation polynomiale aux nœuds équirépartis lorsque le degré augmente.
\end{definition}

\begin{proposition}[Courbure minimale de la spline cubique naturelle]\label{prop:qm:transforms-interpolation-and-algorithmic-differentiation:minimal}
Parmi toutes les fonctions deux fois continûment dérivables passant par les données, la \omterm{def:qm:transforms-interpolation-and-algorithmic-differentiation:spline}{spline cubique naturelle} minimise $\int(f^{\prime\prime})^2$.
\end{proposition}

\begin{proof}
Soit $s$ la spline et $g = f - s$, qui s'annule aux nœuds. En intégrant par parties sur chaque intervalle, $\int s^{\prime\prime}g^{\prime\prime} = [s^{\prime\prime}g']_{x_0}^{x_n} - \int s^{\prime\prime\prime}g' = 0$, car $s^{\prime\prime}$ s'annule aux extrémités et $s^{\prime\prime\prime}$ est
constante sur chaque intervalle, où $\int g' = 0$. D'où $\int(f^{\prime\prime})^2 = \int(s^{\prime\prime})^2 + \int(g^{\prime\prime})^2 \ge \int(s^{\prime\prime})^2$.
\end{proof}

\begin{proposition}[Condition de Fritsch--Carlson]\label{prop:qm:transforms-interpolation-and-algorithmic-differentiation:fc}
Sur un intervalle où les données croissent avec une pente sécante $\delta > 0$, l'interpolante de Hermite cubique de pentes terminales $m_0, m_1 \ge 0$ est monotone si $\alpha = m_0/\delta$ et $\beta = m_1/\delta$ vérifient
$\alpha^2 + \beta^2 \le 9$ (Fritsch et Carlson, 1980).
\end{proposition}

La courbure minimale n'est pas ce que veut un desk de taux. Sur la courbe du Trésor américain du 3 juillet 2023 (rendements à maturité constante d'un mois à trente ans, traités comme des taux zéro pour l'illustration), la
spline des taux zéro produit une courbe à terme qui descend à 2.63~\% à trente ans~; l'interpolation monotone de $rt$, dont la dérivée est le taux à terme, maintient chaque taux à terme entre 3.32~\% et 5.67~\%~;
l'interpolation linéaire des taux zéro donne des taux à terme qui sautent d'au plus 0.43 point de pourcentage aux piliers (voir la \cref{fig:qm:transforms-interpolation-and-algorithmic-differentiation:fwd}). Hagan et West (2006)
passent en revue les méthodes employées pour les courbes. Les polynômes globaux sont pires encore~: sur onze échantillons équirépartis de $1/(1 + 25x^2)$ sur $[-1, 1]$, le polynôme d'interpolation se trompe de 1.92 près des extrémités.

\begin{omfigure}
\begin{tikzpicture}
\begin{axis}[width=11cm,height=5.4cm,xlabel={échéance (années)},ylabel={taux à terme instantané (\%)},
  tick label style={font=\small},label style={font=\small},xmin=0,xmax=30,ymin=2.4,ymax=6,grid=major,grid style={gray!25},
  legend columns=4,legend cell align=left,legend style={font=\footnotesize,at={(0.5,-0.3)},anchor=north,draw=none,fill=none,
  /tikz/every even column/.append style={column sep=6pt}},table/col sep=comma]
\addplot[black!50,thick] table[x=t,y=linear]{figdata/methods/28-transforms-interpolation-and-algorithmic-differentiation/forwards.csv};
\addplot[omThm,very thick] table[x=t,y=spline]{figdata/methods/28-transforms-interpolation-and-algorithmic-differentiation/forwards.csv};
\addplot[omDef,very thick] table[x=t,y=monotone]{figdata/methods/28-transforms-interpolation-and-algorithmic-differentiation/forwards.csv};
\addplot[only marks,mark=*,mark size=1.6pt,black] table[x=t,y=zero]{figdata/methods/28-transforms-interpolation-and-algorithmic-differentiation/pillars.csv};
\legend{taux zéro linéaires,spline cubique naturelle,monotone ($rt$),piliers de taux zéro}
\end{axis}
\end{tikzpicture}

\omcaption{Courbes des taux à terme instantanés implicites des trois interpolations de la courbe du Trésor américain du 3 juillet 2023, avec les rendements des piliers (traités comme des taux zéro). Données~: séries FRED DGS1MO à DGS30 (Conseil des
gouverneurs de la Réserve fédérale, H.15).}\label{fig:qm:transforms-interpolation-and-algorithmic-differentiation:fwd}
\end{omfigure}

\section{Recherche de racines}

\begin{definition}[Méthodes de dichotomie, de la sécante et de Brent]\label{def:qm:transforms-interpolation-and-algorithmic-differentiation:roots}
La \emph{méthode de dichotomie}\index{méthode de dichotomie} divise par deux un encadrement $[a, b]$ tel que $f(a)f(b) < 0$ à chaque étape. La \emph{méthode de la sécante}\index{méthode de la sécante} remplace la dérivée de la \omterm{def:qm:numerical-optimisation-in-practice:newton}{méthode de Newton} par la
pente passant par les deux derniers itérés. La \emph{méthode de Brent}\index{méthode de Brent} (Brent, 1973) conserve un encadrement et effectue une interpolation quadratique inverse ou un pas de sécante lorsque celui-ci reste à l'intérieur de l'encadrement et
le réduit assez vite, et un pas de dichotomie sinon.
\end{definition}

\begin{proposition}[Convergence des méthodes de recherche de racines]\label{prop:qm:transforms-interpolation-and-algorithmic-differentiation:rootconv}
La dichotomie gagne un bit par étape et converge toujours~; la \omterm{def:qm:transforms-interpolation-and-algorithmic-differentiation:roots}{méthode de la sécante} converge localement à l'ordre $(1 + \sqrt5)/2 \approx 1.618$ en racine simple~; la \omterm{def:qm:transforms-interpolation-and-algorithmic-differentiation:roots}{méthode de Brent} converge toujours, jamais beaucoup plus lentement que la dichotomie
(au plus environ le carré du nombre de pas de dichotomie), et superlinéairement près d'une racine simple (Brent, 1973).
\end{proposition}

La volatilité implicite constitue le test de référence. Pour un call à un an ($S = 100$, $r = 3\%$, volatilité vraie 25~\%), la dichotomie sur $[0.001, 5]$ demande 36 pas à chaque strike. À la monnaie, la \omterm{def:qm:numerical-optimisation-in-practice:newton}{méthode de Newton} partant de 20~\% demande 3
itérations, la \omterm{def:qm:transforms-interpolation-and-algorithmic-differentiation:roots}{méthode de la sécante} 4, Brent 5. Pour un strike de 400, le prix vaut $2.5 \times 10^{-7}$ et le vega à 20~\% vaut $8 \times 10^{-9}$~: le premier pas de Newton conduit à une volatilité de 3\,062~\%, où le vega
subit un sous-dépassement, la \omterm{def:qm:transforms-interpolation-and-algorithmic-differentiation:roots}{méthode de la sécante} échoue aussi, tandis que Brent converge en 18. La règle de la maison est une méthode d'encadrement avec un pas intérieur rapide, et Newton seulement à partir d'une estimation dont on sait qu'elle est proche.

\section{Différentiation algorithmique directe et adjointe}

\begin{definition}[Différentiation algorithmique, modes direct et inverse, nombres duaux]\label{def:qm:transforms-interpolation-and-algorithmic-differentiation:ad}
\emph{différentiation algorithmique}\index{différentiation algorithmique} calcule les dérivées d'une fonction donnée sous forme de programme en appliquant la règle de la chaîne à ses opérations élémentaires, exactement à
l'arrondi près. Le \emph{mode direct}\index{mode direct} propage, avec chaque valeur, sa dérivée selon une direction d'entrée~; les \emph{nombres duaux}\index{nombres duaux} $a + b\epsilon$ avec $\epsilon^2 = 0$ le réalisent,
puisque $f(a + b\epsilon) = f(a) + f'(a)b\epsilon$ (Wengert, 1964). Le \emph{mode inverse}\index{mode inverse} enregistre les opérations, puis propage les adjoints $\bar v = \partial y/\partial v$ de la sortie jusqu'à
chaque entrée (Linnainmaa, 1976, issu de son mémoire de master de 1970 sur les erreurs d'arrondi~; Griewank, 2012, recense plusieurs découvertes indépendantes)~; en finance, on l'appelle \emph{mode adjoint}\index{mode adjoint}, ou AAD.
\end{definition}

\begin{definition}[Bande, checkpointing]\label{def:qm:transforms-interpolation-and-algorithmic-differentiation:tape}
Une \emph{bande}\index{bande} est l'enregistrement, réalisé pendant l'évaluation directe, des entrées de chaque opération élémentaire et de ses dérivées partielles locales. Le \emph{checkpointing}\index{checkpointing} ne conserve que certains
états intermédiaires pendant la passe directe et recalcule les opérations entre eux lors de la passe inverse, en échangeant du calcul contre de la mémoire.
\end{definition}

\begin{proposition}[Principe du gradient peu coûteux]\label{prop:qm:transforms-interpolation-and-algorithmic-differentiation:cheap}
Pour une fonction scalaire de $n$ entrées évaluée par $E$ opérations élémentaires, chacune à au plus deux entrées, le \omterm{def:qm:transforms-interpolation-and-algorithmic-differentiation:ad}{mode inverse} calcule le gradient entier avec un nombre d'opérations borné par une petite
constante fois $E$, indépendant de $n$~; le \omterm{def:qm:transforms-interpolation-and-algorithmic-differentiation:ad}{mode direct} et les différences finies demandent $n$ passes (Griewank et Walther, 2008).
\end{proposition}

\begin{proof}
L'enregistrement évalue chaque opération une fois et calcule au plus deux dérivées partielles locales~; le balayage inverse effectue une multiplication suivie d'une addition par couple (opération, entrée) enregistré. Avec $P \le 2E$ couples, le total
vaut $E + 2P \le 5E$ dans ce décompte, quel que soit $n$. Le \omterm{def:qm:transforms-interpolation-and-algorithmic-differentiation:ad}{mode direct} propage une tangente par passe, si bien que le gradient complet demande $n$ passes.
\end{proof}

Le carnet d'exemple de la maison compte 400 piliers de taux zéro sur trente ans, 120 flux trimestriels et 119 caplets évalués par la formule de Black sur les taux forward à trois mois~: 2\,864 opérations pour une valeur. Le balayage inverse
coûte, dans le décompte ci-dessus, 3.58 évaluations~; le bumping coûte 401 et le \omterm{def:qm:transforms-interpolation-and-algorithmic-differentiation:ad}{mode direct} 1\,433 (la \cref{fig:qm:transforms-interpolation-and-algorithmic-differentiation:cost}). Le gradient adjoint concorde
avec le \omterm{def:qm:transforms-interpolation-and-algorithmic-differentiation:ad}{mode direct} à $2 \times 10^{-13}$ près et avec les différences centrales à $9 \times 10^{-6}$ près, sur des sensibilités allant jusqu'à 517 (l'erreur propre aux bumps)~; 213 des 400 sont non nuls, puisque seuls les piliers
voisins d'une date de flux comptent sous interpolation linéaire. Le jumeau C++20, chronométré dans son test d'acceptation, renvoie le gradient d'une fonction semblable à 400 entrées en moins de vingt fois une
évaluation simple (entre cinq et sept dans les exécutions faites pour ce chapitre, allocation de la \omterm{def:qm:transforms-interpolation-and-algorithmic-differentiation:tape}{bande} comprise).

\begin{omfigure}
\begin{tikzpicture}
\begin{loglogaxis}[width=10cm,height=5cm,xlabel={nombre d'entrées $n$},ylabel={coût du gradient (évaluations)},
  tick label style={font=\small},label style={font=\small},xmin=20,xmax=500,ymin=1,ymax=3000,grid=major,grid style={gray!25},
  xtick={25,50,100,200,400},xticklabels={25,50,100,200,400},
  legend columns=3,legend cell align=left,legend style={font=\footnotesize,at={(0.5,-0.3)},anchor=north,draw=none,fill=none,
  /tikz/every even column/.append style={column sep=8pt}},table/col sep=comma]
\addplot[black,very thick,mark=square*,mark size=1.5pt] table[x=n,y=bump]{figdata/methods/28-transforms-interpolation-and-algorithmic-differentiation/adcost.csv};
\addplot[omThm,very thick,mark=*,mark size=1.5pt] table[x=n,y=forward]{figdata/methods/28-transforms-interpolation-and-algorithmic-differentiation/adcost.csv};
\addplot[omDef,very thick,mark=triangle*,mark size=1.8pt] table[x=n,y=reverse]{figdata/methods/28-transforms-interpolation-and-algorithmic-differentiation/adcost.csv};
\legend{bumping ($n + 1$),mode direct,mode inverse}
\end{loglogaxis}
\end{tikzpicture}

\omcaption{Coût du gradient du carnet de taux d'exemple, en évaluations, par nombre d'opérations, en fonction du nombre de piliers de taux zéro~: le bumping et le \omterm{def:qm:transforms-interpolation-and-algorithmic-differentiation:ad}{mode direct} croissent avec $n$, le \omterm{def:qm:transforms-interpolation-and-algorithmic-differentiation:ad}{mode inverse}, lui, non.
Données~: le tutoriel du chapitre.}\label{fig:qm:transforms-interpolation-and-algorithmic-differentiation:cost}
\end{omfigure}

Le prix du \omterm{def:qm:transforms-interpolation-and-algorithmic-differentiation:ad}{mode inverse} est la mémoire~: la \omterm{def:qm:transforms-interpolation-and-algorithmic-differentiation:tape}{bande} conserve chaque opération. Une boucle en pas de temps de 10\,000 pas enregistre 60\,003 nœuds~; avec un checkpoint tous les 100 pas, la passe inverse conserve au plus 201
états (100 checkpoints et un segment) et recalcule la boucle une fois, et le gradient est le même à $10^{-14}$ près. Deux mises en garde pratiques~: un payoff à point anguleux ou à saut a une dérivée nulle ou
indéfinie presque partout, donc l'AAD d'un digital Monte Carlo renvoie zéro et exige un lissage ou un terme de rapport de vraisemblance~; et le code qui branche sur des valeurs n'enregistre que la branche empruntée. Giles et
Glasserman (2006) ont introduit la méthode adjointe dans le calcul des Greeks Monte Carlo en finance.

\section{Tutoriel~: un seul balayage au lieu de 401}\label{tut:qm:transforms-interpolation-and-algorithmic-differentiation:aad}

\begin{tutorial}
\textbf{Objectif.} Calculer un prix par COS à partir d'un exposant de Lévy, construire trois courbes à terme, inverser Black--Scholes avec quatre méthodes de recherche de racines, et différencier un carnet à 400 entrées de trois façons. \textbf{État final~:} les quatre figures
et la concordance des gradients.
\begin{tutsteps}
\item \textbf{La \omterm{def:qm:transforms-interpolation-and-algorithmic-differentiation:tape}{bande}.} Nœuds, balayage inverse et une opération surchargée dans l'en-tête C++20.
\omcode{code/firm/aad/cpp/firm_aad.hpp}{13}{52}{Une bande, son balayage inverse et une multiplication enregistrée (C++20).}
\item \textbf{\omterm{def:qm:transforms-interpolation-and-algorithmic-differentiation:tape}{Checkpointing}.} La référence Python recalcule chaque segment sur une \omterm{def:qm:transforms-interpolation-and-algorithmic-differentiation:tape}{bande} neuve.
\omcode{code/firm/aad/firm_aad.py}{201}{229}{Mode inverse à travers une longue boucle avec checkpoints (Python).}
\item \textbf{Exécuter} \texttt{cos\_convergence()}, \texttt{vg\_check()}, \texttt{gil\_pelaez\_digital()}, \texttt{forward\_curves()}, \texttt{implied\_vol\_iterations()}, \texttt{gradients()} et
\texttt{checkpoint\_demo()} dans \texttt{qm\_transforms.py}~; puis \texttt{fig\_transforms.py}.
\end{tutsteps}
\textbf{Pour aller plus loin.} Remplacer l'interpolation linéaire par une interpolation monotone dans le carnet et observer le nombre de sensibilités non nulles~; différencier le prix COS par rapport aux paramètres de Merton en
exécutant \texttt{cos\_call} sur des \omterm{def:qm:transforms-interpolation-and-algorithmic-differentiation:ad}{nombres duaux}~; ajouter un gamma du second ordre (\omterm{def:qm:transforms-interpolation-and-algorithmic-differentiation:ad}{mode direct} sur \omterm{def:qm:transforms-interpolation-and-algorithmic-differentiation:ad}{mode inverse}).
\end{tutorial}

\section{Construction~: différentiation algorithmique}\label{bld:qm:transforms-interpolation-and-algorithmic-differentiation:aad}

\begin{build}
\bfield{Rôle} Sensibilités exactes de tout code de valorisation au prix de quelques évaluations, pour la bibliothèque de valorisation et le moteur de risque des livres suivants.
\bfield{Interface} Référence Python~: \texttt{Tape}, \texttt{Var}, \texttt{Dual}, \texttt{exp}, \texttt{log}, \texttt{sqrt}, \texttt{ncdf}, \texttt{maximum}, \texttt{gradient}, \texttt{forward\_gradient},
\texttt{bump\_gradient}, \texttt{checkpointed\_gradient}. C++20~: \texttt{firm::aad::Var} sur un \texttt{tape} local à chaque thread, \texttt{Tape::reverse}, \texttt{Dual}. Rust~: \texttt{Var}, \texttt{reverse},
\texttt{clear}, \texttt{Dual}, \texttt{ncdf}.
\bfield{Règles} Chaque gradient AAD est livré avec un contrôle par bump sur un échantillon d'entrées~; la \omterm{def:qm:transforms-interpolation-and-algorithmic-differentiation:tape}{bande} est vidée à chaque valorisation et à chaque thread~; les points anguleux et les sauts sont lissés avant la différentiation~; les longues boucles sont
checkpointées.
\bfield{Tests d'acceptation} \texttt{code/firm/aad/\{tests,cpp,rust\}}~: règles élémentaires contre des dérivées calculées à la main~; nombres d'opérations~; la valeur de la fonction test à 400 entrées et trois composantes du gradient
égales dans les trois langages~; \omterm{def:qm:transforms-interpolation-and-algorithmic-differentiation:ad}{mode direct} contre \omterm{def:qm:transforms-interpolation-and-algorithmic-differentiation:ad}{mode inverse} à $10^{-11}$ près~; bumps centrés et unilatéraux~; \omterm{def:qm:transforms-interpolation-and-algorithmic-differentiation:tape}{checkpointing} égal à la \omterm{def:qm:transforms-interpolation-and-algorithmic-differentiation:tape}{bande} complète~; le rapport de temps C++ inférieur à vingt.
\bfield{Pour aller plus loin} \omterm{def:qm:transforms-interpolation-and-algorithmic-differentiation:ad}{Mode direct} vectoriel~; dérivées secondes~; modèles d'expressions pour réduire la \omterm{def:qm:transforms-interpolation-and-algorithmic-differentiation:tape}{bande}~; adjoints d'une résolution linéaire (chapitre~25) sans enregistrer ses boucles internes.
\end{build}

\begin{omsources}
\item J.~W.~Cooley et J.~W.~Tukey, \emph{Mathematics of Computation} 19, 1965~; P.~Carr et D.~Madan, «~Option valuation using the fast Fourier transform~», \emph{Journal of Computational Finance} 2(4), 1999.
\item F.~Fang et C.~W.~Oosterlee, «~A novel pricing method for European options based on Fourier-cosine series expansions~», \emph{SIAM Journal on Scientific Computing} 31, 2008~; J.~Gil-Pelaez, \emph{Biometrika} 38, 1951.
\item F.~N.~Fritsch et R.~E.~Carlson, «~Monotone piecewise cubic interpolation~», \emph{SIAM Journal on Numerical Analysis} 17, 1980~; P.~S.~Hagan et G.~West, «~Interpolation methods for curve construction~»,
\emph{Applied Mathematical Finance} 13, 2006~; C.~Runge, \emph{Zeitschrift für Mathematik und Physik} 46, 1901.
\item R.~P.~Brent, \emph{Algorithms for Minimization without Derivatives}, Prentice-Hall, 1973.
\item R.~E.~Wengert, \emph{Communications of the ACM} 7, 1964~; S.~Linnainmaa, \emph{BIT} 16, 1976~; A.~Griewank, «~Who invented the reverse mode of differentiation~?~», \emph{Documenta Mathematica}, Extra Volume ISMP, 2012~; A.~Griewank et A.~Walther, \emph{Evaluating Derivatives}, 2\up{e} éd., SIAM, 2008~; M.~Giles et P.~Glasserman,
«~Smoking adjoints~: fast Monte Carlo Greeks~», \emph{Risk}, janvier 2006.
\item Rendements du Trésor~: Board of Governors of the Federal Reserve System, H.15, via FRED (séries DGS1MO à DGS30), consulté le 24 septembre 2026.
\end{omsources}

\section{Exercices}

\begin{exercise}[$\star$]\label{exo:qm:transforms-interpolation-and-algorithmic-differentiation:1}
Calculer le nombre de multiplications complexes qu'exige une \omterm{def:qm:transforms-interpolation-and-algorithmic-differentiation:dft}{transformée de Fourier discrète} directe de longueur 4\,096, et estimer l'ordre de grandeur de celui d'une \omterm{def:qm:transforms-interpolation-and-algorithmic-differentiation:dft}{transformée de Fourier rapide} à base 2.
\end{exercise}

\begin{exercise}[$\star$]\label{exo:qm:transforms-interpolation-and-algorithmic-differentiation:2}
Évaluer $f(x) = x^2e^x$ sur le nombre dual $1 + \epsilon$ et lire $f'(1)$.
\end{exercise}

\begin{exercise}[$\star$]\label{exo:qm:transforms-interpolation-and-algorithmic-differentiation:3}
Calculer le nombre de pas de dichotomie nécessaires pour ramener un encadrement de largeur 5 sous $10^{-10}$.
\end{exercise}

\begin{exercise}[$\star\star$]\label{exo:qm:transforms-interpolation-and-algorithmic-differentiation:4}
Écrire la \omterm{def:qm:transforms-interpolation-and-algorithmic-differentiation:tape}{bande} de $y = \sin(x_1x_2) + x_1$ et effectuer à la main le balayage inverse en $(x_1, x_2) = (1, 2)$.
\end{exercise}

\begin{exercise}[$\star\star$]\label{exo:qm:transforms-interpolation-and-algorithmic-differentiation:5}
Montrer que le forward instantané vaut $f(t) = r(t) + tr'(t)$, et expliquer pourquoi une spline de $r$ de forte pente négative à l'extrémité longue de la courbe produit là de faibles forwards.
\end{exercise}

\begin{exercise}[$\star\star$]\label{exo:qm:transforms-interpolation-and-algorithmic-differentiation:6}
Expliquer pourquoi le delta AAD par trajectoire d'une option digitale calculée par Monte Carlo est nul, et citer deux remèdes.
\end{exercise}

\begin{exercise}[$\star\star\star$]\label{exo:qm:transforms-interpolation-and-algorithmic-differentiation:7}
\emph{Programmation.} Relever les nombres d'itérations nécessaires au calcul de la volatilité implicite pour les quatre méthodes de recherche de racines, aux strikes 50, 70, 100, 150, 250 et 400, et expliquer les échecs.
\end{exercise}

\begin{exercise}[$\star\star\star$]\label{exo:qm:transforms-interpolation-and-algorithmic-differentiation:8}
\emph{Repérer le défaut.} «~Notre gradient AAD diffère de celui obtenu par bumping de $10^{-5}$ sur une sensibilité de 500, donc la bibliothèque AAD comporte une erreur.~»
\end{exercise}

\section{Problème~: un seul balayage au lieu de 401}

\begin{problem}[{Problème du week-end --- 400 sensibilités pour le prix de quelques-unes}]\label{pb:qm:transforms-interpolation-and-algorithmic-differentiation:1}
Un carnet de taux fictif, adossé à une courbe zéro à 400 piliers, détient 120 flux trimestriels et 119 caplets. Le calcul nocturne des risques applique un choc à chaque pilier.

\medskip\noindent\textbf{Partie I --- Prix issus des transformées.}
\begin{enumerate}
\item Comment l'erreur de la \omterm{def:qm:transforms-interpolation-and-algorithmic-differentiation:cos}{méthode COS} évolue-t-elle avec le nombre de termes pour Black--Scholes et Merton~?
\item Que donne la formule de Gil-Pelaez pour le digital de Merton, et par rapport à quoi~?
\item Comment vérifie-t-on le prix variance-gamma calculé par la \omterm{def:qm:transforms-interpolation-and-algorithmic-differentiation:cos}{méthode COS}~?
\item Pourquoi évaluer les calls à partir du développement des puts~?
\item Quand la \omterm{def:qm:transforms-interpolation-and-algorithmic-differentiation:dft}{transformée de Fourier rapide} est-elle l'outil le plus adapté~?
\end{enumerate}

\medskip\noindent\textbf{Partie II --- Courbes et racines.}
\begin{enumerate}[resume]
\item Que deviennent les forwards sous la \omterm{def:qm:transforms-interpolation-and-algorithmic-differentiation:spline}{spline cubique naturelle} des taux zéro, et où~?
\item Que garantit l'interpolation monotone de $rt$, et quelle plage de forwards donne-t-elle~?
\item Quel est le défaut de l'interpolation linéaire des taux zéro~?
\item Comment se comportent les quatre méthodes de recherche de racines à la monnaie et pour un strike de 400~?
\item Quelle est l'erreur de Runge sur onze nœuds équidistants~?
\end{enumerate}

\medskip\noindent\textbf{Partie III --- Le gradient.}
\begin{enumerate}[resume]
\item Combien d'opérations une valorisation effectue-t-elle, et combien de couples (opération, entrée)~?
\item Que coûtent le bumping, le \omterm{def:qm:transforms-interpolation-and-algorithmic-differentiation:ad}{mode direct} et le \omterm{def:qm:transforms-interpolation-and-algorithmic-differentiation:ad}{mode inverse}, en évaluations~?
\item Dans quelle mesure les trois gradients concordent-ils, et pourquoi le bumping diffère-t-il~?
\item Pourquoi seules 213 sensibilités sur 400 sont-elles non nulles~?
\item Qu'apporte le \omterm{def:qm:transforms-interpolation-and-algorithmic-differentiation:tape}{checkpointing} sur une boucle de 10\,000 pas~?
\end{enumerate}

\medskip\noindent\textbf{Partie IV --- Jugement.}
\begin{enumerate}[resume]
\item Quelles sensibilités faut-il encore soumettre au bumping~?
\item Que doit vérifier le moteur de risque chaque nuit~?
\item Qu'ajoute le temps d'exécution en C++ au décompte des opérations~?
\item Énoncer le \emph{résultat nommé}~: le rapport de coût entre le gradient adjoint et une évaluation pour le carnet à 400 entrées, à comparer aux 401 évaluations du bumping, ainsi que l'accord entre les deux.
\item En une phrase~: pourquoi le \omterm{def:qm:transforms-interpolation-and-algorithmic-differentiation:ad}{mode adjoint} est-il le mode naturel pour le risque~?
\end{enumerate}
\end{problem}

\section{Questions d'entretien}

\begin{interviewq}[$\star$ \iqroles{chercheur, développeur}]\label{iq:qm:transforms-interpolation-and-algorithmic-differentiation:1}
Expliquez la différentiation en \omterm{def:qm:transforms-interpolation-and-algorithmic-differentiation:ad}{mode direct} et en \omterm{def:qm:transforms-interpolation-and-algorithmic-differentiation:ad}{mode inverse}, et dites dans quel cas chacune est la moins coûteuse.
\end{interviewq}

\begin{interviewq}[$\star\star$ \iqroles{développeur}]\label{iq:qm:transforms-interpolation-and-algorithmic-differentiation:2}
Comment implémenteriez-vous le \omterm{def:qm:transforms-interpolation-and-algorithmic-differentiation:ad}{mode inverse} en C++, et qu'est-ce qui en limite l'usage en production~?
\end{interviewq}

\begin{interviewq}[$\star\star$ \iqroles{chercheur}]\label{iq:qm:transforms-interpolation-and-algorithmic-differentiation:3}
Comment calculeriez-vous de façon robuste une volatilité implicite pour un strike quelconque~?
\end{interviewq}

\begin{interviewq}[$\star\star$ \iqroles{chercheur}]\label{iq:qm:transforms-interpolation-and-algorithmic-differentiation:4}
Quelle interpolation utiliseriez-vous pour une courbe d'actualisation, et pourquoi pas une \omterm{def:qm:transforms-interpolation-and-algorithmic-differentiation:spline}{spline cubique naturelle} des taux zéro~?
\end{interviewq}

\begin{interviewq}[$\star\star$ \iqroles{chercheur}]\label{iq:qm:transforms-interpolation-and-algorithmic-differentiation:5}
Comment évalueriez-vous une option européenne dans un modèle connu seulement par sa \omterm{def:qm:probability-at-speed:charfn}{fonction caractéristique}~?
\end{interviewq}

\begin{interviewq}[$\star\star\star$ \iqroles{développeur, risque}]\label{iq:qm:transforms-interpolation-and-algorithmic-differentiation:6}
Votre Monte Carlo en AAD donne un delta nul pour une option à barrière. Que s'est-il passé et que faites-vous~?
\end{interviewq}
